\begin{figure}[t]
\centering
\begin{tikzpicture}[
    layer/.style={
        draw, rounded corners=4pt, minimum width=11cm,
        minimum height=0.85cm, font=\small, align=center,
        line width=0.7pt
    },
]
% ---- Probabilistic Execution Plane (background) ----
\fill[classical!12, rounded corners=8pt]
    (-6.0,-0.6) rectangle (6.0,3.5);

% ---- Deterministic Control Plane (background) ----
\fill[modelnative!12, rounded corners=8pt]
    (-6.0,3.9) rectangle (6.0,6.8);

% ---- Plane labels (in dedicated header space above topmost layer) ----
\node[font=\small\bfseries, text=classical!80!black, anchor=west,
      fill=classical!8, fill opacity=0.6, text opacity=1,
      rounded corners=2pt, inner sep=2pt] at (-5.8,3.2)
    {概率执行平面};
\node[font=\small\bfseries, text=modelnative!80!black, anchor=west,
      fill=modelnative!8, fill opacity=0.6, text opacity=1,
      rounded corners=2pt, inner sep=2pt] at (-5.8,6.5)
    {确定性控制平面};

% ---- L1 ----
\node[layer, fill=classical!25, draw=classical!70]
    (L1) at (0,0) {L1\enspace 物理执行层：矩阵运算与采样};

% ---- L2 ----
\node[layer, fill=classical!25, draw=classical!70]
    (L2) at (0,1.2) {L2\enspace 推理服务层：KV缓存\,·\,连续批处理\,·\,解码调度};

% ---- L3 ----
\node[layer, fill=classical!25, draw=classical!70]
    (L3) at (0,2.4) {L3\enspace 上下文管理层：热/温/冷分层\,·\,虚拟上下文};

% ---- Interface line ----
\draw[dashed, thick, black!40] (-6.0,3.7) -- (6.0,3.7);

% ---- L4 ----
\node[layer, fill=modelnative!25, draw=modelnative!70]
    (L4) at (0,4.5) {L4\enspace 语义接口层：工具schema\,·\,能力标注\,·\,审计};

% ---- L5 ----
\node[layer, fill=modelnative!25, draw=modelnative!70]
    (L5) at (0,5.7) {L5\enspace 编排层：任务规划\,·\,权限审批\,·\,失败恢复};

% ---- L6 (above both planes) ----
\node[layer, fill=gray!20, draw=gray!60]
    (L6) at (0,7.4) {L6\enspace 应用层：用户任务\,·\,Agent应用\,·\,治理审计};

% ---- Vertical arrows ----
\draw[-{Latex[length=2mm]}, thick, black!35] (L1) -- (L2);
\draw[-{Latex[length=2mm]}, thick, black!35] (L2) -- (L3);
\draw[-{Latex[length=2mm]}, thick, black!35] (L3) -- (L4);
\draw[-{Latex[length=2mm]}, thick, black!35] (L4) -- (L5);
\draw[-{Latex[length=2mm]}, thick, black!35] (L5) -- (L6);

% ---- Governance bar ----
\draw[thick, gray!50] (-6.0,8.0) -- (6.0,8.0);
\node[font=\footnotesize, text=gray!70] at (0,8.3) {治理覆盖全栈（公理5\,·\,公理6）};

\end{tikzpicture}
\caption{双平面架构：概率执行平面（蓝，L1--L3）与确定性控制平面（橙，L4--L5）的分离}
\label{fig_dual_plane}
\end{figure}
